\chapter{The Jacobian as an abelian variety}
\label{chap:Jacobian}

Let \(C\) be a smooth proper geometrically irreducible curve over a field
\(k\), with genus \(g(C)\).  The Jacobian is constructed uniformly as the
identity component \(\Pic^0_{C/k}\) of the Picard scheme.  It is an abelian
variety of dimension \(g(C)\); when \(g(C)=0\), it is \(\Spec k\).

\section{Albanese objects}

\begin{definition}\leanok
[Albanese universal property of a pointed curve]
  \label{def:IsAlbanese}
  \lean{AlgebraicGeometry.IsAlbanese}
  Let \(P \in C(k)\).  An abelian variety \(J\) together with a pointed
  morphism \(\iota_P \colon C \to J\) is an \emph{Albanese object} for
  \((C,P)\) if every pointed morphism \(f \colon C \to A\) to an abelian
  variety factors uniquely through \(\iota_P\): there is a unique morphism of
  \(k\)-schemes \(g \colon J \to A\) such that
  \[
    P \circ \iota_P = \eta_J,
    \qquad
    f = \iota_P \circ g.
  \]
\end{definition}

\begin{remark}
  \label{rem:IsAlbanese_typeclasses}
  The phrase ``abelian variety'' includes the group structure, properness,
  smoothness, and geometric irreducibility of both \(J\) and every target
  \(A\).  These are hypotheses of the universal property, rather than extra
  conclusions inside it.
\end{remark}

\subsection{The universal morphism}

\begin{definition}\leanok
[The Albanese universal morphism]
  \label{def:IsAlbanese_ofCurve}
  \uses{def:IsAlbanese}
  \lean{AlgebraicGeometry.IsAlbanese.ofCurve}
  An Albanese structure on \((C,P,J)\) determines a distinguished morphism
  \(\iota_P \colon C \to J\), called its universal morphism.
\end{definition}

\begin{lemma}\leanok
[Pointed condition]
  \label{lem:IsAlbanese_comp_ofCurve}
  \uses{def:IsAlbanese, def:IsAlbanese_ofCurve}
  \lean{AlgebraicGeometry.IsAlbanese.comp_ofCurve}
  The universal morphism satisfies \(P \circ \iota_P = \eta_J\).
\end{lemma}

\begin{lemma}\leanok
[Universal factorisation]
  \label{lem:IsAlbanese_exists_unique_ofCurve_comp}
  \uses{def:IsAlbanese, def:IsAlbanese_ofCurve}
  \lean{AlgebraicGeometry.IsAlbanese.exists_unique_ofCurve_comp}
  Every pointed morphism \(f \colon C \to A\) to an abelian variety factors
  uniquely as \(f=\iota_P\circ g\) through a morphism of \(k\)-schemes
  \(g \colon J \to A\).
\end{lemma}

\begin{theorem}\leanok
[Uniqueness of the Albanese object]
  \label{thm:IsAlbanese_unique}
  \uses{def:IsAlbanese, def:IsAlbanese_ofCurve}
  \lean{AlgebraicGeometry.IsAlbanese.unique}
  If \((J_1,\iota_1)\) and \((J_2,\iota_2)\) are Albanese objects for
  \((C,P)\), there is a unique morphism \(e \colon J_1\to J_2\) satisfying
  \(\iota_2=\iota_1\circ e\).
\end{theorem}

\begin{proof}\leanok
  The universal property of \(J_1\) applied to \(\iota_2\) gives
  \(e \colon J_1\to J_2\); that of \(J_2\) applied to \(\iota_1\) gives
  \(e' \colon J_2\to J_1\).  Both \(e\circ e'\) and the identity of \(J_1\)
  factor \(\iota_1\) through itself, so uniqueness gives
  \(e\circ e'=\mathrm{id}_{J_1}\); similarly
  \(e'\circ e=\mathrm{id}_{J_2}\).  The original universal property gives the
  required uniqueness of \(e\).
\end{proof}

\begin{remark}
  \label{rem:IsAlbanese_unique_iso}
  Consequently the Albanese object is unique up to a unique isomorphism
  compatible with the universal morphism.
\end{remark}

\section{A uniform Jacobian witness}

\begin{definition}\leanok
[Jacobian witness]
  \label{def:JacobianWitness}
  \uses{def:IsAlbanese, def:genus}
  \lean{AlgebraicGeometry.JacobianWitness}
  A \emph{Jacobian witness} for \(C\) consists of a single \(k\)-scheme \(J\),
  an abelian-variety structure on \(J\), a proof that \(J\) is smooth of
  relative dimension \(g(C)\), and, for every \(P\in C(k)\), an Albanese
  structure on \((C,P,J)\).
\end{definition}

\begin{remark}
  \label{rem:JacobianWitness_quantifier_order}
  The same \(J\) works for every marked point.  Replacing \(P\) by another
  rational point translates the universal morphism by an automorphism of
  \(J\), so it does not change the underlying abelian variety.
\end{remark}

\begin{remark}
  \label{rem:JacobianWitness_smooth_redundancy}
  Smoothness is implied by smoothness of relative dimension \(g(C)\).  It is
  nevertheless useful to record both the abelian-variety property and the
  dimension assertion as parts of the witness.
\end{remark}

\subsection{The Picard witness \(J=\Pic^0_{C/k}\)}
\label{sec:picardJacobianWitness}

\begin{theorem}\leanok
[The Picard identity component is a Jacobian witness]
  \label{def:picardJacobianWitness}
  \uses{def:JacobianWitness, def:genus, thm:fga_pic_representability,
    thm:pic0_isAbelianVariety, thm:pic_zero_dimension_equals_genus,
    thm:albanese_universal_property_northstar}
  \lean{AlgebraicGeometry.picardJacobianWitness}
  The identity component \(\Pic^0_{C/k}\) carries a Jacobian witness for
  \(C\).
\end{theorem}

\begin{proof}
  Representability of the relative Picard functor gives \(\Pic_{C/k}\).
  Its identity component is a group scheme; the structure theorem for the
  identity component of the Picard scheme makes it smooth, proper, and
  geometrically irreducible.  The tangent-space identification
  \[
    T_0\Pic^0_{C/k}\cong H^1(C,\mathcal O_C)
  \]
  gives relative dimension \(g(C)\).  Finally, the Abel--Jacobi morphism and
  its universal property give an Albanese structure for each \(P\in C(k)\).
  These data form the required witness.
\end{proof}

\begin{theorem}\leanok
[Existence of a Jacobian witness]
  \label{thm:nonempty_jacobianWitness}
  \uses{def:JacobianWitness, def:picardJacobianWitness}
  \lean{AlgebraicGeometry.nonempty_jacobianWitness}
  Every smooth proper geometrically irreducible curve has a Jacobian witness.
\end{theorem}

\begin{proof}
  Take the witness on \(\Pic^0_{C/k}\) supplied by
  Theorem~\ref{def:picardJacobianWitness}.
\end{proof}

\begin{definition}\leanok
[Chosen Jacobian witness]
  \label{def:jacobianWitness}
  \uses{def:JacobianWitness, thm:nonempty_jacobianWitness}
  \lean{AlgebraicGeometry.jacobianWitness}
  Fix one Jacobian witness for \(C\).
\end{definition}

\section{The Jacobian and its structure}

\begin{definition}\leanok
[Jacobian]
  \label{def:Jacobian}
  \uses{def:genus, def:JacobianWitness, def:jacobianWitness}
  \lean{AlgebraicGeometry.Jacobian}
  The \emph{Jacobian} \(\Jac(C)\) is the underlying \(k\)-scheme of the
  chosen Jacobian witness.  Equivalently, \(\Jac(C)=\Pic^0_{C/k}\) up to the
  unique isomorphism of Albanese objects.
\end{definition}

\begin{theorem}\leanok
[Group structure]
  \label{thm:Jacobian_grpObj}
  \uses{def:Jacobian, def:jacobianWitness}
  \lean{AlgebraicGeometry.Jacobian.instGrpObj}
  The Jacobian is a group scheme over \(k\).
\end{theorem}

\begin{proof}
  This is the group structure carried by the chosen Jacobian witness.
\end{proof}

\begin{theorem}\leanok
[Dimension and smoothness]
  \label{thm:Jacobian_smooth_genus}
  \uses{def:Jacobian, def:genus, def:jacobianWitness}
  \lean{AlgebraicGeometry.Jacobian.smoothOfRelativeDimension_genus}
  The structural morphism \(\Jac(C)\to\Spec k\) is smooth of relative
  dimension \(g(C)\).
\end{theorem}

\begin{proof}
  This is the dimension assertion carried by the chosen Jacobian witness.
\end{proof}

\begin{theorem}\leanok
[Properness]
  \label{thm:Jacobian_proper}
  \uses{def:Jacobian, def:jacobianWitness}
  \lean{AlgebraicGeometry.Jacobian.instIsProper}
  The structural morphism \(\Jac(C)\to\Spec k\) is proper.
\end{theorem}

\begin{proof}
  This is the properness assertion carried by the chosen Jacobian witness.
\end{proof}

\begin{theorem}\leanok
[Geometric irreducibility]
  \label{thm:Jacobian_geomIrred}
  \uses{def:Jacobian, def:jacobianWitness}
  \lean{AlgebraicGeometry.Jacobian.instGeometricallyIrreducible}
  The structural morphism \(\Jac(C)\to\Spec k\) is geometrically irreducible.
\end{theorem}

\begin{proof}
  This is the geometric-irreducibility assertion carried by the chosen
  Jacobian witness.
\end{proof}

\begin{remark}
  If \(g(C)=0\), then \(\Jac(C)=\Spec k\).  If \(g(C)=1\) and \(C(k)\) is
  nonempty, the choice of a rational point identifies \(C\) with its
  Jacobian.  For \(g(C)\geq2\), the Jacobian has dimension \(g(C)\), whereas
  \(C\) has dimension one.
\end{remark}

\section{The Albanese property of \(\Pic^0\)}
\label{sec:Jacobian_routeA4_albaneseUP}

\begin{theorem}
[Albanese universal property of \(\Pic^0_{C/k}\)]
  \label{thm:albanese_universal_property_northstar}
  \uses{lem:abel_jacobi_morphism, lem:symmetric_product_av_map,
    lem:symmetric_product_to_jacobian, lem:descent_through_birational_sigma,
    lem:albanese_eq_iff_symmetricPower_eq, thm:rational_map_to_av_extends,
    thm:rigidity_lemma, lem:av_regular_map_is_hom,
    lem:hom_additivity_over_product}
  \lean{AlgebraicGeometry.Scheme.Pic.albaneseUP}
  Fix \(P\in C(k)\), put \(J=\Pic^0_{C/k}\), and let
  \(f^P(Q)=[\mathcal O_C(Q-P)]\).  For every pointed morphism
  \(\varphi\colon C\to A\) to an abelian variety there is a unique
  group-scheme homomorphism \(\psi\colon J\to A\) such that
  \(\varphi=f^P\circ\psi\).
\end{theorem}

\begin{proof}
  For \(g=g(C)>0\), the morphism
  \[
    (Q_1,\ldots,Q_g)\longmapsto
      \varphi(Q_1)+\cdots+\varphi(Q_g)
  \]
  is symmetric and therefore descends to \(C^{(g)}\).  The Abel map
  \(C^{(g)}\to J\) is birational onto \(J\), so this descent defines a
  rational map \(J\dashrightarrow A\).  A rational map from a nonsingular
  variety to an abelian variety extends uniquely to a regular morphism;
  denote the extension by \(\psi\).  By construction
  \(\varphi=f^P\circ\psi\), and \(\psi(0)=0\), so rigidity makes \(\psi\) a
  homomorphism.  If \(\psi'\) is another such homomorphism, then \(\psi\) and
  \(\psi'\) agree on the sum of \(g\) copies of \(f^P(C)\), which is all of
  \(J\); hence \(\psi=\psi'\).  When \(g=0\), \(J=\Spec k\) and the result
  follows from rigidity.
\end{proof}
